\documentclass[11pt]{article}
\usepackage[margin=1in]{geometry}
\usepackage{hyperref}
\usepackage{enumitem}
\usepackage{booktabs}
\usepackage{xcolor}
\title{ProbOS --- Month 2 Plan}
\author{Nisong Monyimba}
\date{\today}

\begin{document}
\maketitle

\section*{Theme}
Sequential inference plus bringing the C++ kernel into the Python package.
Month 1 built a forward-simulation kernel: draw priors, propagate,
summarise, explain. Month 2's central gap is \textbf{inference} --- given
observed data, update beliefs about latent state and parameters online.
That capability is the particle filter, and it is the spine around which
every other Month 2 workstream is organised.

\section*{Week 5 --- Particle Filter Core}
\textbf{Goal:} a working sequential Monte Carlo filter, validated against
a known analytical case before touching \texttt{BatteryModel2Cell}.

\begin{itemize}[leftmargin=*]
  \item \texttt{python/src/particle\_filter.py}
  \begin{itemize}
    \item \texttt{ParticleFilter} class: init from prior, \texttt{predict()}
      step (reuses \texttt{Model.forward\_batch()}), \texttt{update()} step
      (likelihood weighting), \texttt{resample()} step (systematic
      resampling)
    \item Effective sample size (ESS) diagnostic, resample-when-ESS-low
      trigger
  \end{itemize}
  \item Validate against a linear-Gaussian model with a known Kalman
    filter solution first, then apply to \texttt{BatteryModel2Cell} with
    synthetic noisy temperature observations.
  \item Study guide: Chopin \& Papaspiliopoulos (2020) Ch 8--10 on
    resampling schemes; Naesseth et al. (2019) for the unifying SMC
    algorithm template.
\end{itemize}
\textbf{Exit test:} filter posterior mean tracks true latent state within
theoretical bounds on the linear-Gaussian case; ESS never collapses
without triggering resampling.

\section*{Week 6 --- pybind11: C++ Kernel Enters the Python Package}
\textbf{Goal:} stop maintaining two parallel implementations. Bind
\texttt{BatteryCell} and \texttt{MonteCarloEngineOMP} into Python via
pybind11.

\begin{itemize}[leftmargin=*]
  \item \texttt{cpp/bindings/} --- pybind11 module definition
  \item \texttt{setup.py} / \texttt{CMakeLists.txt} updates for building
    the extension
  \item Cross-validation test: C++ and Python engines must agree on
    percentiles for the same seed within floating-point tolerance
\end{itemize}
\textbf{Exit test:} \texttt{import probos\_cpp} works from a plain
\texttt{pytest} run in CI.

\section*{Week 7 --- FastAPI Service Layer}
\textbf{Goal:} expose the kernel over HTTP.

\begin{itemize}[leftmargin=*]
  \item \texttt{POST /simulate} --- run \texttt{MonteCarloEngine}, return
    percentiles and convergence summary
  \item \texttt{POST /sensitivity} --- run \texttt{SobolSensitivity},
    return S1/ST indices
  \item \texttt{POST /filter} --- run \texttt{ParticleFilter} against
    posted observation data
  \item Pydantic schemas mirroring the \texttt{Model}/\texttt{Distribution}
    constructor signatures
\end{itemize}
\textbf{Exit test:} a \texttt{curl} POST to \texttt{/simulate} returns
valid JSON matching a direct Python call to the same
\texttt{MonteCarloEngine}.

\section*{Week 8 --- Cross-Discipline Batch 2 + Month 2 Retrospective}
\textbf{Goal:} show the particle filter also generalises across domains.

\begin{itemize}[leftmargin=*]
  \item Port the ED queue model to a filtering setting (infer true
    arrival rate from observed queue-length time series)
  \item Port the clinical trial model to sequential monitoring, using the
    particle filter instead of batch nested-MC
  \item \texttt{docs/retrospectives/month2\_retrospective.md} plus
    \texttt{month3\_plan.md}
\end{itemize}
\textbf{Exit test:} both filtering examples pass validation against
ground truth, with test coverage matching Month 1's standard.

\section*{Standing Month 2 commitments (every week)}
\begin{itemize}[leftmargin=*]
  \item Run \texttt{check\_ci.sh} after every push, no exceptions
  \item Update \texttt{docs/study/study\_guide.md} same-session as any
    new file needing external reading
  \item mypy strict and ruff clean before any commit
  \item No claimed speedup or result without a genuine benchmark backing
    it
\end{itemize}


\section*{What was actually accomplished (retrospective)}

\begin{itemize}[leftmargin=*]
  \item \textbf{Week 5 -- ParticleFilter Core}: bootstrap SIR particle
    filter, validated against an exact closed-form 1D Kalman filter
    (implemented directly, not imported) on a linear-Gaussian test
    case. Genuine mathematical proof, not a smoke test: posterior
    mean/std converge to the exact solution, with Monte Carlo error
    provably shrinking as $N$ grows from 100 to 2000.
  \item \textbf{Week 6 -- pybind11 bindings}: the Week 4 C++ kernel
    (\texttt{BatteryCell}, \texttt{MonteCarloEngineOMP}) bound into
    the Python package as the \texttt{probos\_cpp} extension, ending
    the two-disconnected-implementations problem. Single-step
    numerical correctness proven to \texttt{rtol}$\le$1e-8 against the
    Python reference. A known P05--P95 spread discrepancy between the
    two engines (simplified C++ parameter sampling vs full Python
    priors) was documented and explicitly tested-for, not silently
    accepted.
  \item \textbf{Week 7 -- FastAPI Service Layer}: \texttt{/simulate},
    \texttt{/sensitivity}, \texttt{/filter} endpoints, each verified
    to match a direct Python kernel call at \texttt{rtol}=1e-10.
    Resource-exhaustion bounds enforced via Pydantic before any
    request reaches kernel code.
  \item \textbf{Week 8 -- Cross-Discipline Batch 2}: \texttt{ParticleFilter}
    demonstrated across two more domains (hospital operations,
    clinical trials), each validated against an independently
    computable ground truth. See
    \texttt{docs/monthly\_plans/month2/week8/main.tex} for full detail.
  \item \textbf{Comprehensive mid-month hardening pass} (between Weeks
    7 and 8, prompted by explicit user request to leave ``no stone
    unturned'' before continuing): a systematic, file-by-file audit
    of the entire repository found and fixed a substantial number of
    real issues that had accumulated silently, including:
    \begin{itemize}
      \item \texttt{mypy --strict} had never actually been run against
        \texttt{python/tests/} in this project's history -- running it
        surfaced $\sim$150 findings across 3 files (all fixture-typing
        or deliberate-negative-test patterns, zero real bugs), fixed
        properly and the CI/audit scope permanently extended to cover
        it going forward.
      \item The full ruff configuration declared in
        \texttt{pyproject.toml} (\texttt{I}, \texttt{UP}, \texttt{B},
        \texttt{SIM} rule sets) had never actually been enforced --
        every invocation throughout the project used a narrower
        \texttt{--select E,F,W} override. Running the full config
        surfaced 46 findings, including a genuinely tautological test
        assertion that was rewritten to test its own stated intent
        rather than papered over.
      \item A genuine license mismatch: \texttt{README.md} said MIT
        throughout while the actual \texttt{LICENSE} file is
        Apache-2.0 -- a real, standalone factual/legal error, fixed
        outright.
      \item Eight vestigial \texttt{sys.path.insert()} hacks removed
        across \texttt{python/src/} and \texttt{python/examples/},
        predating the project's editable install and verified
        genuinely unnecessary via real subprocess import tests (not
        just AST parsing), both from the repo root and from outside
        the repo directory.
      \item \texttt{Model.validate\_params()}/\texttt{validate\_state()}
        had docstrings falsely claiming to be called automatically by
        \texttt{MonteCarloEngine} -- genuinely wired into
        \texttt{MonteCarloEngine}, \texttt{ParticleFilter}, and
        \texttt{SobolSensitivity} for real, adding a genuine
        early-failure safety check across all three engines.
      \item Several smaller off-by-one documentation errors (the
        \texttt{Distribution} ABC has 5 concrete classes, not 4 as
        README claimed; the server registers 4 example models, not 5
        as its own docstring claimed) and 3 stale manuscript figure
        duplicates causing the compiled PDF to embed outdated plots.
    \end{itemize}
    This pass is documented in detail across the individual commits
    from that period; it is summarised here because it materially
    affected Month 2's overall quality, not just Week 8's.
\end{itemize}

\section*{What went right}
\begin{enumerate}[leftmargin=*]
  \item Every new capability (particle filter, pybind11 bindings,
    FastAPI endpoints, Week 8's cross-discipline filters) was
    validated against an independently computable ground truth before
    being trusted -- the same discipline established in Month 1,
    consistently maintained.
  \item Real bugs were found and fixed rather than hidden or
    rationalised at every stage this month: the prior-mismatch bug in
    Week 8's clinical trial validation, the reproducibility bug in the
    ED queue ground-truth generator, the false claims in
    \texttt{validate\_params}'s docstring, and the license mismatch
    are all documented honestly in their respective commits.
  \item The mid-month comprehensive audit demonstrates the project's
    own quality-assurance tooling (\texttt{pre\_week\_audit.sh},
    \texttt{docs/standards/quality\_standards.md}) is not just
    aspirational -- when genuinely applied comprehensively (not just
    the narrower subset previously enforced), it found real,
    substantive issues, and those enforcement gaps were then
    permanently closed rather than re-opening next week.
\end{enumerate}

\section*{What was harder than expected}
\begin{enumerate}[leftmargin=*]
  \item Designing Week 8's clinical trial filter required real
    domain reasoning, not mechanical reuse: the naive approach
    (directly reusing \texttt{ClinicalTrialModel.forward\_batch()}'s
    own stochastic simulation inside the standard predict/update loop)
    would have caused severe weight degeneracy from comparing
    independently-simulated discrete outcomes against real data. The
    correct design (deterministic state transitions against known
    real data, informative likelihood scoring via particle
    parameters) took real first-principles thinking to arrive at.
  \item The mid-month audit revealed that several ``we already
    checked this'' assumptions were false -- notably, that the full
    declared ruff/mypy configuration was actually being enforced
    anywhere. This is a useful, humbling lesson: a checklist that is
    declared but not actually run provides zero real assurance,
    regardless of how thorough it looks on paper.
\end{enumerate}

\section*{Numbers at Month 2 close}
\begin{itemize}[leftmargin=*]
  \item Python tests: 368/368 passing
  \item C++ tests: 13/13 passing
  \item mypy strict: 0 errors (entire \texttt{python/} tree, including
    \texttt{python/tests/} -- a scope extension made this month)
  \item ruff: 0 warnings (full declared configuration, not a narrowed
    subset -- a scope extension made this month)
  \item Coverage: 90.80\%
  \item 0 bandit findings, 0 known dependency CVEs
\end{itemize}
\end{document}
